\section{Defences and Detection}
\label{sec:defences}

\subsection{What an operator can do}

\paragraph{Give each agent its own execution role.}
This is the one with the largest effect and the clearest evidence behind it:
\nSharedRoleRuntimes of \nMultiAgentRuntimes multi-agent deployments do not do
it today. A shared role makes the takeover unnecessary, because the agents
already hold the same permissions, and it removes the account's ability to
say which agent made a call. Separating the roles does not by itself stop the
takeover, but it is what makes every other control in this list possible.

\paragraph{Name the role instead of matching its name.}
The gate in the vendor's managed policy is a resource pattern over the role
name, so two agents whose roles both match \iam{*BedrockAgentCore*} are not
separated by it at all. A per agent policy granting \iam{iam:PassRole} on the
exact role ARN the agent needs, in place of
\iam{BedrockAgentCoreFullAccess}, names the role instead of matching a
pattern.

This is worth stating as a defence because the vendor's own policy checker
does not treat the pattern as a problem. IAM Access Analyzer raises
\texttt{PASS\_ROLE\_WITH\_STAR\_IN\_RESOURCE} on \iam{iam:PassRole} with a
wildcard resource, and its reference documentation records an exception for
this exact shape: a resource such as \iam{arn:aws:iam::*:role/Service*} may
produce no finding, because some services put their own namespace in the names
of the roles they create~\cite{awsAccessAnalyzerChecks}. The convention that
makes the pattern look specific is the same convention that places every
agent's execution role inside it, so an operator running the vendor's checker
over the vendor's policy is told nothing.

\paragraph{Narrow the trust policy.}
The documented trust policy for an execution role scopes
\iam{sts:AssumeRole} with \iam{aws:SourceArn} set to a wildcard over the whole
account. Replacing the wildcard with the specific runtime ARN closes the
target side of the edge, and the role then cannot be reused by
another runtime.

\paragraph{Separate the identity that deploys from the identity that runs.}
\api{UpdateAgentRuntime} is a deployment operation. An agent's own execution
role rarely needs it, and where it does the grant should name the runtime it
may update. The corpus contains \nBlockedAgentToAgent candidates that are one
agent reaching another and are held back only by \iam{iam:PassRole}; in those
deployments the update permission is already present on the wrong identity.

\paragraph{Constrain the command route as well.}
The controls above govern who may replace an agent's artifact. They do nothing
about the permission that runs a command inside an existing runtime, which
reaches the same role more cheaply and has a service control policy of its
own~\cite{garcia2026agentcorepaths}. Apply both.

\subsection{Whether the first defence refuses the cell that succeeded}
\label{sec:arn-repair}

The defence above is a proposal until something evaluates it, so it is
evaluated here with the engine that decides the real call. The
\iam{iam:PassRole} statement is read from the live managed policy rather than
retyped, a second copy is made with only the \iam{Resource} element replaced
by the ARN of the one role an agent is entitled to pass, and both are put to
\api{SimulateCustomPolicy} against three targets with
\iam{iam:PassedToService} set to the service. Nothing is created and no
runtime is involved; the artifact is
\texttt{evaluation/\allowbreak benchmark/\allowbreak arn\_condition\_repair.json}
and the script that writes it reads the policy afresh on each run.

Table~\ref{tab:repair} gives the \nRepairDecisions decisions. The row that
matters is the second: the shipped statement is \repairShippedSibling{} to
pass a sibling agent's role, because the console named that role inside the
pattern, and the same statement pinned to an ARN is \repairPinnedSibling. The
last row is the one that makes it a repair and not a removal, since the
pinned statement still allows the agent to pass its own role, which is what
the permission is for.

\begin{table}[t]
\centering
\caption{The first defence, put to the IAM policy simulator. The shipped
statement and the same statement with its resource pinned to one role ARN,
each against three targets. The repair denies the cell that succeeded in
Section~\ref{sec:case-study} and leaves the agent's own pass working.}
\label{tab:repair}
\small
\begin{tabular}{@{}llL{22mm}@{}}
\toprule
\textbf{Role being passed} & \textbf{Shipped v21} & \textbf{Pinned to one ARN} \\
\midrule
the agent's own role & allowed & allowed \\
a sibling the console named & allowed & denied \\
a sibling named anything else & denied & denied \\
\bottomrule
\end{tabular}
\end{table}

Two limits belong with this. The simulator decides authorisation, so what is
measured is the decision and not the \api{UpdateAgentRuntime} call being
refused end to end; the live cells in Section~\ref{sec:case-study} are what
show the call itself. And the repair moves the problem rather than dissolving
it: an account now has to get one policy right per agent, where before it
attached one policy everywhere, which is work that falls on whoever writes the
deployment.

\subsection{What an operator cannot fix this way}

Scoping an agent's data permissions does not help: agent~A never touches
agent~B's data, because the code it installs runs as agent~B and uses
agent~B's own permissions. A deployment where every role is perfectly scoped to its own agent is still
exposed to it. Controls at the model
layer do not help either, because the cloud permissions downstream are the
same whatever ran the attacker's code.

\subsection{What the analysis hands an operator}

A rule based scanner returns a verdict: the deployment violates the rule or it
does not. On this corpus that verdict is uninformative, because the resource
type is not in the rule set and all \nCheckovFailed findings are about
something else. But the deeper limitation is the shape of the answer rather
than the coverage. A deployment that is one grant away from a cross-agent
takeover and a deployment that is nowhere near one both come back clean, and
nothing in the output separates them.

Evaluating each delegation rule as a conjunction changes the shape of the
answer. When a rule does not fire, the analysis knows which conjunct failed,
so the output is not ``no path'' but ``no path, and this is the term that is
missing''. Over this corpus that term is the same one every time: all
\nBlocked candidates lack \iam{iam:PassRole} on the target role, and
Algorithm~\ref{alg:reach} carries the witness, so each one names the file and
the principal a reader can open.

Three things follow for an operator, and they are properties of the output
format, not new measurements. A blocked candidate is a ranking, because
a deployment one permission from a takeover is a different risk from one that
is several. It is also a specific instruction, since the permission not to
grant is named rather than inferred. And it is a regression test: the
\nPassRoleWide configuration where that permission is already granted broadly
is the same finding with the distance at zero, which is how a deployment
crosses from one class to the other when somebody widens a policy. The
numbers in Section~\ref{sec:results} are a census of the corpus; the same
output on one deployment is a work list.

\subsection{Detection}

Two signals follow directly from the mechanism, and both are available in
ordinary cloud audit logs.

The first is an \api{UpdateAgentRuntime} call whose \api{roleArn} differs from
the runtime's previous value. A deployment pipeline changes the artifact often
and the role almost never, so a change of role on an existing runtime is rare
enough to alert on rather than to sample.

The second is an \api{UpdateAgentRuntime} call made by a principal that is
itself an agent execution role. Agents are not usually deployers. Neither
signal is specific to an attack; both are specific enough to review.

Both miss the name condition cell of Section~\ref{sec:case-study}, where the
role does not change and the attacker is an ordinary principal. The third
signal is a comparison, not an event: reading each runtime's artifact
back and checking it against the object the pipeline published names the
runtimes whose code is no longer what was deployed. It is the only one of the
three that still fires once the attacker is gone.

\subsection{What the platform could change}

If replacing an artifact and changing an execution role were separate
operations with separate permissions, the takeover path would need two grants
instead of one, and the common case of updating an agent that keeps its
identity would stop requiring \iam{iam:PassRole} at all. Short of that, a
condition key expressing ownership instead of a name pattern would let the
managed policy say what it appears to say.
